\chapter{ИЛЛЮСТРАТИВНЫЙ РАЗБОР МОДЕЛИ}
\label{app:example}

Рассмотрим описание: «Чашка кофе имеет начальную температуру 180 °F и остывает до комнатной температуры 70 °F с характерным временем 20 минут». Это иллюстративный пример представления, а не утверждение о независимо полученном результате LLM.

Краткий brief содержит один stock (Teacup Temperature), внешнюю room temperature, потерю тепла и отрицательную обратную связь. В канонической форме модель записывается следующим образом:

\begin{align}
    T(0)&=180,\\
    R(t)&=70,\\
    \tau&=20,\\
    q(t)&=\frac{T(t)-R(t)}{\tau},\\
    \frac{dT}{dt}&=-q(t).
\end{align}

Здесь $T$ и $R$ имеют температурные единицы, $\tau$ --- минуты, а $q$ --- температурные единицы в минуту. Поток направлен от stock к окружающей среде, поэтому производная отрицательна. Если exporter получит $dT/dt=q$, topology check отклонит модель, поскольку поток объявлен исходящим.

Аналитическое решение:

\begin{equation}
    T(t)=R+(T(0)-R)e^{-t/\tau}=70+110e^{-t/20}.
\end{equation}

При $t=30$ ожидается приблизительно $87{,}85$ температурных единиц. Численный Euler не обязан совпадать с этим значением точно: ошибка зависит от шага. Поэтому в benchmark задаётся допуск, а не требование побайтового равенства с аналитической формулой.

Пример показывает, почему недостаточно получить правдоподобный CLD. Для рабочей модели необходимы начальное значение, размерность $\tau$, правильное направление потока и численно корректная производная. Эти свойства могут быть проверены автоматически до запуска симуляции.
